\newglossaryentry{test}
{
    name={test},
    description={test},
    plural={tests}
}
\newglossaryentry{ggt}
{
    type=dinge,
    name={ggt},
    description={test},
    plural={tests}
}
\newglossaryentry{paarweise disjunkt}
{
    name={paarweise disjunkt},
    description={Ein Mengensystem $M$, also eine Menge von Mengen, nennt man \gls{paarweise disjunkt}, wenn jeweils zwei verschiedene Mengen $A, B \in M$ disjunkt sind},
    plural={paarweise disjunkte}
}
\newglossaryentry{multimenge}
{
    name={Multimenge},
    description={Eine Multimenge ist eine Sammlung von Elementen, bei der die
Wiederholung von Elementen erlaubt ist. Anders als in einer normalen Menge, in der
jedes Element nur einmal vorkommt, kann jedes Element in einer Multimenge beliebig
oft auftreten},
    plural={Multimengen}
}
\newglossaryentry{partition}
{
    type=dinge,
    name={$\mathcal{P}$},
    description={test},
    plural={\mathcal{P}}
}
\newglossaryentry{k}
{
    type=dinge,
    name={$k$},
    description={test},
    plural={k}
}
\newglossaryentry{r}
{
    type=dinge,
    name={$r$},
    description={test},
    plural={r}
}
\newglossaryentry{b}
{
    type=dinge,
    name={$b$},
    description={test},
    plural={b}
}
\newglossaryentry{urbild}
{
    type=dinge,
    name={$f^{-1}$},
    description={Urbild},
    plural={-1},
}
\newglossaryentry{Homomorphismus}{
name={Homomorphismus},
description={Abbildung zwischen zwei Gruppen, die die Struktur erhält}
}
\newglossaryentry{multinomialkoeffizient}
{
    type=dinge,
    name={$\binom{n}{r_1, r_2, \dots, r_k}$},
    description={test},
    plural={\binom{n}{r_1, r_2, \dots, r_k}}
}
\newglossaryentry{kSnk}
{
    type=dinge,
    name={\( k! \cdot S(n, k) \)},
    description={Anzahl der surjektiven Abbildungen \( X \to Y \), falls \( |X| = n \) und \( |Y| = k \). Auch: Anzahl der Möglichkeiten \( n \) Objekte in \( k \) Fächer zu verteilen, sodass kein Fach leer bleibt.}
}
\newglossaryentry{Sn}
{
    type=dinge,
    name={$S_n$},
    description={Menge der Permutationen von \( \{1, \dots, n\} \)},
    plural={S_n}
}
\newglossaryentry{S}
{
    type=dinge,
    name={$S(n,k)$},
    description={Die Anzahl der Partitionen einer $n$-elementigen Menge in $k$ Teile - Stirlingzahlen zweiter Art},
    plural={S}
}
\newglossaryentry{snk}
{
    type=dinge,
    name={$s(n,k)$},
    description={Die Anzahl der Permutationen einer $n$-elementigen Menge, die genau $k$ Zyklen haben - Stirlingzahlen erster Art},
    plural={s}
}
\newglossaryentry{AI}
{
    type=dinge,
    name={$A_I$},
    description={Durchschnitt aller Mengen $A_i$ mit $i$ in der Menge $I$: $\bigcap_{i \in I} A_i$},
    plural={A_I}
}
\newglossaryentry{chiA}
{
    type=dinge,
    name={$\chi_A$},
    description={},
    plural={\chi}
}

\newglossaryentry{zarn}
{
    type=dinge,
    name={$\mathbb{Z} / 2 \mathbb{Z}[x]\backslash (x^n + 1)$},
    description={test},
    plural={$(2)$}
}

\newglossaryentry{zavn}
{
    type=dinge,
    plural={$(1)$},
    description={test},
    name={$(\mathbb{Z} / 2 \mathbb{Z}[x])^n$}
}

\newglossaryentry{ideal}
{
    name={Ideal},
    description={test},
    plural={Ideale}
}

\newglossaryentry{Körper}
{
    name={Körper},
    description={Ein Körper ist in der Algebra eine ausgezeichnete algebraische Struktur, in der die Addition, Subtraktion, Multiplikation und Division auf eine bestimmte Weise durchgeführt werden können},
    plural={Körper}
}
\newglossaryentry{Koeffizient}
{
    name={Koeffizient},
    description={Generell ist damit eine feste Zahl (Konstante) gemeint, die du mit einer variablen Zahl, wie zum Beispiel $x$ multiplizierst: $3x$. In diesem Beispiel ist $3$ der Koeffizient },
    plural={Koeffizienten}
} 
\newglossaryentry{Polynom}
{
    name={Polynom},
    description={Summe von Vielfachen von Potenzen einer Variablen bzw. Unbestimmten},
    plural={Polynome}
}
\newglossaryentry{Ring}
{
    name={Ring},
    description={Ein Ring ist eine Menge RR, die zwei Operationen hat: Addition und Multiplikation, mit folgenden Eigenschaften: \begin{itemize}
        \item Addition
        \item Multiplikation
        \item evtl. multiplikatives neutrales Element
        \item evtl. Nullteilerfreiheit
    \end{itemize}},
    plural={Ringe}
}
\newglossaryentry{Faktorisierung}
{
    name={Faktorisierung},
    description={Zerlegung eines Objekts in mehrere nichttriviale Faktoren},
    plural={Faktorisierungen}
}
\newglossaryentry{Restklassenring}
{
    name={Restklassenring},
    description={test},
    plural={Restklassenringe}
} 
\newglossaryentry{L̈osungsraum}
{
    name={L̈osungsraum},
    description={test},
    plural={L̈osungsäauem}
}
\newglossaryentry{Kontrollmatrix}
{
    name={Kontrollmatrix},
    description={test},
    plural={tests}
}
\newglossaryentry{nullteilerfrei}
{
    name={nullteilerfrei},
    description={eine Struktur (Ring,..), die keine nichttrivialen \gls{Nullteiler} enthält.},
    plural={nullteilerfrei}
}
\newglossaryentry{Nullteiler}
{
    name={Nullteiler},
    description={Element $a$ eines Rings, für das es ein vom Nullelement 0 verschiedenes Element $b$ gibt, so dass $\frac{a}{b} = 0$},
    plural={Nullteiler}
}

\newglossaryentry{Isomorphismus}
{
    name={Isomorphismus},
    description={Eine Abbildung zwischen zwei mathematischen Strukturen, durch die Teile einer Struktur auf bedeutungsgleiche Teile einer anderen Struktur umkehrbar eindeutig (bijektiv) abgebildet werden},
    plural={Isomorphismen}
}